\chapter{Conclusiones}

El proyecto del **Sistema de Gestión Integral Universitario**, actualmente en sus fases iniciales de planificación y definición de requerimientos, se perfila como una solución estratégica para la Universidad Nacional. Tras la etapa de análisis preliminar y la consolidación de la Especificación de Requerimientos de Software (SRS), se pueden establecer las siguientes conclusiones sobre su potencial y el camino a seguir:

\begin{itemize}
    \item \textbf{Viabilidad y Pertinencia del Proyecto:} Se ha confirmado la viabilidad técnica y operativa del Sistema de Gestión Integral Universitario, el cual responde a necesidades críticas identificadas en la gestión de evidencias para acreditación SINAES y la administración de tiempos de jornada universitaria. Este proyecto es pertinente y fundamental para la modernización y optimización de procesos académicos y administrativos clave de la UNA.

    \item \textbf{Claridad en la Definición de Requerimientos:} A través de un proceso riguroso de recopilación y consolidación, se ha logrado definir una Especificación de Requerimientos de Software (SRS) completa, coherente y suficientemente detallada para guiar las próximas fases de diseño y desarrollo. Los requerimientos funcionales y no funcionales de ambos módulos han sido claros y validados con los usuarios clave.

    \item \textbf{Alineación con Estándares y Metodologías Modernas:} La adopción de metodologías ágiles como Scrum y el uso de estándares de desarrollo y seguridad modernos (React/Next.js, NestJS, MongoDB, Prisma, Google Drive API, JWT, AES-256) proveen una base tecnológica y de gestión robusta para el proyecto, prometiendo un desarrollo eficiente y de alta calidad.

    \item \textbf{Potencial de Impacto Significativo:} Se proyecta que el sistema contribuirá a la reducción de tiempos administrativos (ej. clasificación de evidencias), la eliminación de duplicidades, la mejora en la precisión de la planificación de recursos humanos y el fortalecimiento de la trazabilidad y seguridad de la información institucional. Estos beneficios son cruciales para los procesos de acreditación y la eficiencia operativa de la UNA.

    \item \textbf{Base para Futuras Fases:} Los requerimientos definidos y la arquitectura modular propuesta para el Sistema de Gestión Integral Universitario sientan una base sólida para su desarrollo en el curso de Ingeniería en Sistemas I y para futuras expansiones y mejoras por parte de la Universidad, consolidando una plataforma integral de gestión de la calidad.
\end{itemize}

\section{Logros hasta la Fecha (Fase de Planificación)}

En esta etapa inicial del proyecto, los principales logros se centran en la preparación y el establecimiento de una base sólida:

\begin{itemize}
    \item \textbf{Conformación y Coordinación del Equipo de Proyecto:} El equipo de estudiantes ha demostrado capacidad de organización y colaboración efectiva, estableciendo canales de comunicación claros con la profesora y los clientes.

    \item \textbf{Definición Exhaustiva del Alcance y Requerimientos:} Se han identificado y documentado los requerimientos funcionales y no funcionales esenciales para ambos módulos, culminando en un SRS consolidado que servirá como guía principal para el desarrollo.

    \item \textbf{Planificación Detallada del Desarrollo:} Se ha establecido un cronograma de Sprints y una estrategia de desarrollo ágil que permitirá la entrega incremental de funcionalidades, optimizando el uso del tiempo disponible en el curso.

    \item \textbf{Análisis de Viabilidad Confirmado:} Los estudios de factibilidad técnica, operativa y económica han ratificado que el proyecto es viable con los recursos y tecnologías disponibles en el entorno universitario.
\end{itemize}

En síntesis, aunque el desarrollo del software apenas comienza, la fase de planificación ha sido exitosa en establecer una dirección clara, un conjunto de requerimientos bien definidos y una base estratégica sólida para el Sistema de Gestión Integral Universitario.